% year: תשפ"ד
% semester: ב
% moed: ג
% number: 1ב
% categories: func-limits
% source: exams/מבחנים/תשפד/מועד ג תשפד סמסטר ב.pdf

\begin{question}
חשבו
\[ \lim_{x\longrightarrow-8}\frac{\sqrt{1+x}-3}{2+\sqrt[3]{x}} \]
\end{question}

\begin{hint}
בדקו קודם את תחום ההגדרה של $\sqrt{1+x}$ ליד $x=-8$.
\end{hint}

\begin{solution}
\textbf{כפי שהשאלה מודפסת}, הביטוי $\sqrt{1+x}$ מוגדר (בממשיים) רק עבור $x\ge-1$, ולכן הפונקציה אינה מוגדרת באף סביבה מנוקבת של $x=-8$, והגבול \textbf{אינו מוגדר}. (גם ההצבה $x=-8$ נותנת $\sqrt{-7}$.)

ככל הנראה נפלה טעות דפוס, והכוונה ל-$\frac{\sqrt{1-x}-3}{2+\sqrt[3]{x}}$ — אז בנקודה $x=-8$ המונה $\sqrt9-3=0$ והמכנה $2+(-2)=0$, גבול מהצורה $\frac00$. נפתור גבול זה. נכפיל ב"צמודים": במונה $\sqrt{1-x}-3=\frac{(1-x)-9}{\sqrt{1-x}+3}=\frac{-(x+8)}{\sqrt{1-x}+3}$, ובמכנה, לפי $a^3+b^3=(a+b)(a^2-ab+b^2)$ עם $a=2,\ b=\sqrt[3]x$:
\[ 2+\sqrt[3]{x}=\frac{8+x}{4-2\sqrt[3]{x}+\sqrt[3]{x^2}}. \]
לכן עבור $x\neq-8$:
\[ \frac{\sqrt{1-x}-3}{2+\sqrt[3]{x}}=-\frac{4-2\sqrt[3]{x}+\sqrt[3]{x^2}}{\sqrt{1-x}+3}\xrightarrow[x\to-8]{}-\frac{4+4+4}{3+3}=-2, \]
לפי רציפות הביטוי האחרון ב-$x=-8$. (אותה תוצאה מתקבלת מכלל לופיטל: $\frac{-\frac{1}{2\sqrt{1-x}}}{\frac{1}{3\sqrt[3]{x^2}}}\to\frac{-1/6}{1/12}=-2$.)
התשובה: כפי שמודפס — הגבול אינו מוגדר; בגרסה המתוקנת $\sqrt{1-x}$ — הגבול הוא $\boxed{-2}$.
\end{solution}
